\subsection{Replacing a component removes an old carrying amount}
\label{sec:ppe-replacement}

An asset can remain in use while one of its components is replaced. The company
then owns a new component and has given up the remaining benefits of the old
one. Recording only the new purchase would preserve an asset that no longer
exists. The register therefore needs component identity, original cost,
accumulated depreciation, remaining life and the dates of replacement and use.
Paragraphs 13 and 70 of the retained AASB 116 require the replaced part's
carrying amount to be removed when an eligible replacement is recognised
\citep{aasb116dated}.

Consider a synthetic company with cash and contributed capital of 200 on
1 January 2024. It pays 100 for an administrative asset: a shell costing 80
with an eight-year useful life, and a significant component costing 20 with a
two-year life. Both are available immediately, have zero residual values and
consume benefits evenly. The company admits the cost model. These are
authored facts, including recognition eligibility and the absence of any
additional impairment; real assets require evidence for those assessments.
Administrative use makes depreciation an expense. A production asset might
instead contribute to inventory cost under paragraph 49.

Acquisition debits the two asset costs by 80 and 20 and credits cash by 100.
Paragraphs 43--44 require separate depreciation of significant components;
the full year's charge is $80/8+20/2=20$. Debit depreciation expense 20 and
credit each component's accumulated depreciation by 10. At 31 December 2024,
cash is 100, net PPE is 80 and equity is 180 after the loss of 20. No revenue
or other operating activity is included in this asset-only example.

On 1 January 2025 the company replaces the old component. Its remaining
carrying amount is $20-10=10$. Debit accumulated depreciation 10 and disposal
loss 10, and credit the old component's cost 20. The replacement costs 30,
is paid in cash and is available immediately for three years of even use,
again with zero residual. Its entry debits new PPE 30 and credits cash 30.
Cash is now 70 and PPE is $70+30=100$; equity is 170. Retaining the old
component would instead report PPE of 110 and miss the loss of 10, although
both acquisition entries would balance.

The requirement to remove the old part also applies when it was not
depreciated separately. Paragraph 70 permits replacement cost as an indication
of the old part's cost at acquisition if its carrying amount cannot practicably
be determined. That does not make today's invoice the old part's remaining
value. Acquisition-date cost and consumption would still need a supported
estimate. Ordinary servicing has a different consequence: paragraph 12
expenses it as incurred. A separately assumed maintenance payment of 4 would
debit maintenance expense and credit cash by 4. Capitalising it would balance
but overstate both PPE and profit by 4 under those facts.

During 2025 the shell and new component each depreciate by 10. Immediately
before sale on 31 December, gross active cost is $80+30=110$, accumulated
depreciation is $20+10=30$, and net PPE is 80. Cash remains 70 and equity is
150. The buyer then obtains control and pays 90, with no selling costs or
sale-and-leaseback arrangement. Remove gross cost 110 and accumulated
depreciation 30, debit cash 90, and credit disposal gain 10. Paragraphs
67--72 locate the disposal at control transfer and measure the gain as net
proceeds less carrying amount. The gain is not ordinary revenue; routine
sales of formerly rented assets have a separate rule at 68A.

\begin{center}
\begin{tabular}{@{}lrr@{}}
\toprule
Pretax amount & 2024 & 2025 including sale\\
\midrule
Depreciation expense & 20 & 20\\
Replacement loss & 0 & 10\\
Disposal gain & 0 & 10\\
Profit & $-20$ & $-20$\\
Closing cash / net PPE & 100 / 80 & 160 / 0\\
Closing equity & 180 & 160\\
Operating cash & 0 & 0\\
Investing cash & $-100$ & 60\\
\bottomrule
\end{tabular}
\end{center}

The cash-flow statement separates the 2025 purchase of 30 from the sale
proceeds of 90. Under the retained AASB 107.16(a)--(b), both are investing
flows. The indirect operating reconciliation removes depreciation and the
non-operating loss and gain from profit:
\[
  \operatorname{CFO}_{2024}=-20+20=0,\qquad
  \operatorname{CFO}_{2025}=-20+20+10-10=0.
\]
Lifetime investing cash is $-100-30+90=-40$, equal to the equity reduction
once no assets remain. An ordinary servicing payment would instead be
operating cash; paragraph 16 does not permit investing classification merely
because an expense relates to equipment \citep{aasb107dated}.

These cash-flow comparisons concern calendar 2024 and 2025. The retained
AASB 107 compilation is dated 31 December 2023 and its official application
record covers annual starts from 1 January 2024 to before 1 January 2026.
The retained AASB 116 compilation is dated 31 December 2022 and covers annual
starts from 1 January 2023 to before 1 January 2027. Their different intervals
must remain attached to the respective rules. Tax has not been computed in
these pretax statements; no tax-law or complete reporting claim follows.

\subsection{Revaluation needs each asset's past gains and losses}
\label{sec:ppe-revaluation}

Revaluation changes the amounts to be depreciated and can change profit and
other comprehensive income (OCI) differently. A class total alone cannot
determine the split. For the Australian for-profit case here, paragraphs
39--40 refer to the earlier decreases and surplus of the same asset. The
Australian not-for-profit exceptions permit different class-level treatment;
they cannot be imported into this case \citep{aasb116dated}.

A second synthetic company owns two administrative machines, A and B, which
form the entire class for which it elects the revaluation model. They became
available on 1 January 2023 at costs 100 and 40, with useful lives of five
and four years, zero residuals and even use. Cash after their acquisition was
60, funded by historical capital of 200. Every fair value used below is an
explicitly assumed reliable measurement. Paragraphs 29, 31 and 36--38 require
the class policy, sufficiently regular measurements and coverage of the whole
class; they do not infer missing fair values from the ledger. No additional
impairment is assumed after the separate assessment required by paragraph 63.

The 2023 history establishes the opening for the 2024 statement comparison.
After depreciation of 20 and 10, A and B carry 80 and 30. At 31 December 2023
their assumed fair values are 120 and 24. A increases by 40, recorded in OCI
and accumulated in its revaluation surplus. B decreases by 6; with no surplus
of its own, that decrease is expense in profit or loss. The net class increase
of 34 therefore cannot all be OCI. Such an entry would omit B's expense of
6 and understate OCI by the same amount while leaving total equity unchanged.

For clarity, this company chooses paragraph 35(b)'s treatment of accumulated
depreciation. At each revaluation it first debits accumulated depreciation and
credits gross PPE by the accumulated amount, leaving the net amount as the
adjusted gross balance. It then posts the revaluation difference. The prior
gross cost, accumulated charge and original-cost history remain in the asset
register. This choice changes neither the net revaluation nor its allocation
between profit and OCI.

The opening 2024 balances are consequently cash 60, PPE $120+24=144$,
capital 200, retained earnings $-36$ from depreciation and B's loss, and
surplus 40. Assets and equity both equal 204. The register also remembers
B's unreversed profit-or-loss decrease of 6. This prior-year derivation uses
the retained AASB 116 scope; it does not assert a 2023 cash-flow presentation
under the later AASB 107 compilation.

Let $L_i$ denote an asset's outstanding earlier revaluation decrease in
profit or loss, and $S_i$ its remaining surplus credit. For an increase
$u_i>0$, the part reversing that earlier loss goes to profit; the rest goes
to OCI. For a decrease $d_i>0$, existing surplus absorbs the reduction through
OCI before any excess reaches profit. The operative amounts are
\[
\begin{aligned}
 \text{profit gain}_i&=\min(u_i,L_i),&
 \text{OCI increase}_i&=u_i-\min(u_i,L_i),\\
 \text{OCI decrease}_i&=\min(d_i,S_i),&
 \text{profit loss}_i&=d_i-\min(d_i,S_i).
\end{aligned}
\]
An increase reduces $L_i$ by its profit gain and adds its OCI increase to
$S_i$. A decrease reduces $S_i$ by its OCI decrease and adds its profit loss
to $L_i$. These are revaluation-history rules; the separate impairment
reversal ceiling is considered in the impairment method. No surplus is
transferred during use in this example. If that permitted transfer were
elected, the later surplus available to absorb a decrease would be smaller.

For 2024, the revalued amounts are depreciated over the remaining four and
three years. The charges are $120/4=30$ and $24/3=8$, leaving carrying
amounts 90 and 16. The company then revalues the whole class to 45 and 22.
A falls by 45: its surplus absorbs 40 through OCI and the remaining 5 is
a profit-or-loss expense. B rises by 6, reversing its own earlier loss
through profit. After eliminating accumulated depreciation, gross PPE is
45 and 22, with surplus zero. Profit is $-38-5+6=-37$ and OCI is $-40$.
Cash 60 plus PPE 67 equals closing equity 127. A now retains a prior
revaluation loss of 5; B retains none.

In 2025, the remaining three-year life of A and two-year life of B give
depreciation $45/3=15$ and $22/2=11$. Carrying amounts become 30 and 11.
At 31 December their assumed fair values are 50 and 11. A's increase of
20 first reverses its loss of 5 through profit, then adds 15 to OCI and
surplus. B's valuation is included even though its monetary adjustment is
zero. Before sale, cash 60 and PPE 61 equal equity 121.

Immediately after that measurement, the buyer obtains both machines and pays
their same-date values, 50 for A and 11 for B. Debit cash 61 and remove PPE
61; the disposal gain is zero. The company elects the transfer permitted by paragraph 41:
debit A's surplus 15 and credit retained earnings 15 directly. This transfer
creates no income, new OCI or cash. Routing it through profit would falsely
raise 2025 profit from $-21$ to $-6$.

\begin{center}
\begin{tabular}{@{}lrr@{}}
\toprule
Pretax amount & 2024 & 2025 including sale\\
\midrule
Depreciation & 38 & 26\\
Revaluation loss / gain in profit & 5 / 6 & 0 / 5\\
Disposal gain & 0 & 0\\
Profit / OCI & $-37$ / $-40$ & $-21$ / 15\\
Closing cash / net PPE & 60 / 67 & 121 / 0\\
Closing retained earnings / surplus & $-73$ / 0 & $-79$ / 0\\
Closing equity & 127 & 121\\
Operating cash & 0 & 0\\
Investing cash & 0 & 61\\
\bottomrule
\end{tabular}
\end{center}

The 2024 equity movement is $204-37-40=127$. In 2025, profit of $-21$
and OCI of 15 leave total equity at 121. Retained earnings move from $-73$
to $-79$: the year's loss reduces them by 21 and the direct transfer increases
them by 15. That internal transfer does not change total equity. Operating
cash follows independently from the year's profit and noncash effects:
\[
 \operatorname{CFO}_{2024}=-37+38+5-6=0,\qquad
 \operatorname{CFO}_{2025}=-21+26-5=0.
\]
The entire 61 sale proceeds appear in investing cash. Tax consequences still
need legal tax bases, applicable rates and the profit/OCI origin of each
difference under AASB 112, as AASB 116.42 requires.

AccDSL must therefore preserve the class and its members, each valuation's
date and evidence, original cost and accumulated charges, individual surplus
and loss history, component exits and the chosen surplus-transfer policy.
An absent valuation leaves that class remeasurement unresolved; it is not
permission to select only appreciating assets. A missing disposal date or
old component value likewise prevents the affected calculation. The journals
and these economic changes belong to the same authorised transition. The
zero adjustment for B illustrates why a valuation record may change even
when no monetary line is needed. Together with the state-only advance-allocation
rule above, this preserves the information needed
by the next depreciation, revaluation or disposal.
